\section{The Propers: Themes and Movement}
\zlabel{triptych:concise:themes:start}
The saints belong to Christ through death, and the people celebrating them
ask Christ to complete their own conversion. The feast celebrates the
holiness he has given while teaching the living to seek the mercy, purity,
and peace they still need.

\subsection{Rejoicing and asking mercy}
The Introit calls the Church to rejoice in the Lord and names the Son as
the one whom the angels praise. Its Psalm verse says that praise befits
the upright. Augustine understands uprightness as submission to God:
joy in the Lord requires a will willing to accept him. Schuster describes
the saints as members of Christ sanctified in every condition of life,
including laborers and slaves. Honoring them gives thanks for holiness
received from God.\footnote{Augustine, Ps.~32, enarratio I, \S1
(NPNF, Psalm XXXIII, \S1); Schuster, \work{Sacramentary}, V,
pp.~209--210.}

The Collect asks for mercy through the saints' many intercessions. The
Sunday commemoration adds a petition for release from sins contracted
through frailty. These are the prayers of people who belong to God and
still need forgiveness. The celebration on 1 November 2026 has these
additional Sunday prayers, while retaining the feast's own Lesson,
Gospel, chants, and the Common Preface. The Creed is appointed.

\subsection{The people John hears counted and sees worshipping}
Before the uncountable multitude appears, John hears an exact number:
144,000, with twelve thousand from each named tribe. Bede explains the
number as the Church's completeness. He appeals to membership in Abraham's
seed through faith in Christ and to the apostles' worldwide mission.
He also permits an alternative account of the sequence: Israel receives
the Gospel first, followed by the Gentiles. Both explanations include the
nations; they identify the two companies differently.\footnote{Bede,
\work{Explanation of the Apocalypse}, I, on Apoc.~7:4,9,
pp.~46, 52; Gal.~3:29.}

The multitude's white robes and palms signify baptism and victory through
Christ's Cross in Bede's exposition. Its acclamation ascribes salvation
to God and the Lamb. Bede understands the closing praise as acknowledgment
of gifts received. This supplies one reason for the living to praise the
Son while honoring the saints. Read beside the Alleluia's invitation,
the robes also suggest that conversion begins with a gift: the labor
of discipleship follows admission to Christ's people.\footnote{Bede,
\work{Explanation}, I, on Apoc.~7:9--12, pp.~52--53.}

\subsection{Seeking good and receiving rest}
The Gradual promises that those seeking God will lack no good; the Gospel
blesses hunger for justice and persecution. Augustine answers the devout
poor person's difficulty by naming goods of the heart: justice, truth,
charity, faith, and endurance. He then cites the hunger Beatitude and
identifies Christ as living bread. The promise concerns what a person
seeks in God, even when material poverty remains.\footnote{Augustine,
Ps.~33, enarratio II, \S\S14--15 (NPNF, Psalm XXXIV,
\S\S13--14).}

The Alleluia appoints Christ's invitation to the burdened. Chrysostom
explains it through the following verses about Christ's meek and humble
yoke. He expressly connects humility with the beginning of the Beatitudes.
Humility relieves pride's anxiety, and obedience frees the disciple from
sin's heavier bondage. His account also retains toil, the cross, and
future glory. The refreshment Christ gives can therefore be present while
the Gospel's costly fidelity continues.\footnote{Matt.~11:25--30;
Chrysostom, \work{Matthew}, XXXVIII.3--5.}

\subsection{Offering with the saints and learning their peace}
The Offertory proclaims the peace of righteous souls who appear to have
perished. Augustine invokes Wisdom's description of souls in God's hand
while distinguishing their security from their bodies' awaited
resurrection. In \work{The City of God}, he also reasons from altar
remembrance that the departed faithful remain Christ's members. Their
communion helps explain why the living can honor them and seek their
prayers while still asking salvation for themselves.\footnote{Augustine,
Ps.~87, \S5 (NPNF, Psalm LXXXVIII, \S4);
\work{City of God}, XX.9.}

Augustine's distinction concerns the difference between Christ and his
members. Christ freely lays down his life and takes it up; the righteous
dead await the resurrection already accomplished in their Head. Their
souls' peace is compatible with the body's death and future restoration.
The persecutor has killed someone whose soul remains in God's care.
Schuster explains the chant through that failed claim of
possession: the martyr belongs to God.\footnote{Augustine, Ps.~87,
\S5; Schuster, \work{Sacramentary}, V, p.~211.}

Gregory quotes Wisdom's account of apparent death while explaining a
righteous person's despised preaching and killed body. Through the
Spirit, the example awakens virtue in others. His following explanation
concerns the sufferer's own bodily resurrection. These are different
results of fidelity: conversion among the living and restoration of
the dead.\footnote{Gregory, \work{Moralia}, XII.5.6--7.}

The feast's Secret asks that the offerings honor all the just and benefit
the offerers through God's mercy. Thanksgiving in the Common Preface
is through Christ; human voices seek admission among the angels. Present
worship already addresses the God whom the saints praise, while those
praying still seek the help needed to persevere.

Bede permits the angels in John's vision to represent the worshipping
people, or to be heavenly spirits rejoicing at Gentile salvation.
His latter explanation can be compared with the Introit's angels and
the Preface's heavenly powers. The human congregation asks admission
among those praises. Its thanksgiving and the multitude's acclamation
both acknowledge salvation received from God and the Lamb, while
the congregation continues to ask for mercy.\footnote{Bede,
\work{Explanation}, I, on Apoc.~7:11--12, p.~53.}

The Sunday Secret's request is for growth in service. Schuster understands
it through ministry at the altar, where the sacrifice of praise is offered
through Christ. The worshippers ask God to complete what he has already
given them without their merit. In the first interpretation, their service
is offered within the Church whose departed members remain Christ's own.
In the second, that increase is part of the conversion for which they pray.
The two prayers concern the benefit needed by the living as they honor
the just.\footnote{Schuster, \work{Sacramentary}, III, p.~188;
MR 1962, pp.~417, 703, Secrets.}

The Communion repeats purity, peacemaking, and persecution for justice.
The promise of seeing God requires a pure heart because God is seen
inwardly. Augustine's comparison is with eyes able to receive light;
purification prepares the person for the promised sight. His account of
peace also explains the person's dependence: reason governs desire while
itself submitting to Christ. Chrysostom names the outward work of bringing
divided people together and connects it with the Son's reconciliation
of the alienated. Persecution then tests fidelity to justice. The Gospel's
final address qualifies reviling further: it must be false and for Christ.
Peace and fidelity can require the disciple to withstand hostility while
continuing to seek reconciliation.

The final prayers ask lasting joy, the saints' supplication,
and protection amid danger. Schuster understands that protection as grace
to stand within trial. The living seek both the saints' company and the
conversion which the repeated Beatitudes describe.\footnote{Augustine,
\work{Sermon on the Mount}, I.2.8--I.3.10; Chrysostom,
\work{Matthew}, XV.7; Schuster, \work{Sacramentary}, III,
p.~189; V, p.~211.}
\zlabel{triptych:concise:themes:end}
\clearpage
